% Image credits for the photographs used in this volume (Book 2 only).
% Machine-readable license records live in images/book2/CREDITS.md (and
% images/book1/CREDITS.md for the three photographs reused from Book 1);
% keep the files in sync when adding or removing a photograph.
\chapter*{चित्र-श्रेय}

% Under bidi=basic a Latin run is ONE directional box, so a long author or
% institution name cannot be broken at all; \allowbreak inside it does nothing.
% The Arabic edition reported this page as the book's only overfull boxes.
% \sloppy must be in force when the paragraph is BROKEN, hence the \par before
% \endgroup -- without it the group closes first and the tolerance is restored.
\begingroup\sloppy

इस पुस्तक के रेखाचित्र मौलिक हैं। छायाचित्र अपने-अपने मुक्त अनुज्ञापत्रों के
अंतर्गत प्रयुक्त हैं; उनके रचयिताओं के प्रति आभार:

\begin{itemize}
  \item \emph{Gregor Mendel} (अर्धसूत्री विभाजन का अध्याय): छायाचित्र,
        उन्नीसवीं शताब्दी, Wikimedia Commons, सार्वजनिक अधिकार-क्षेत्र।
  \item \emph{Alexander Fleming} (प्रतिजैविक प्रतिरोध का अध्याय):
        छायाचित्र, Imperial War Museums, 1943, Wikimedia Commons के
        माध्यम से, सार्वजनिक अधिकार-क्षेत्र।
  \item \emph{लसीकाणु से मुकुलित होता एचआईवी} (उपार्जित प्रतिरक्षा का
        अध्याय): C.~Goldsmith का इलेक्ट्रॉन सूक्ष्मचित्र, Centers for
        Disease Control and Prevention, Wikimedia Commons के माध्यम से,
        सार्वजनिक अधिकार-क्षेत्र।
  \item \emph{डार्विन की फिंचें} (विविधीकरण का अध्याय): John Gould का
        उत्कीर्णन, 1845, Wikimedia Commons, सार्वजनिक अधिकार-क्षेत्र।
  \item \emph{संवेदी होमंक्युलस} (गति का अध्याय): OpenStax
        \emph{Anatomy and Physiology} से चित्रांकन, Wikimedia Commons
        के माध्यम से, CC~BY~3.0.
  \item \emph{Charles Darwin} (जातिवृत्तीय वृक्षों का अध्याय):
        Henry Maull और John Fox का छायाचित्र, लगभग 1854, Wikimedia
        Commons, सार्वजनिक अधिकार-क्षेत्र।
  \item \emph{Edward Jenner} (उपार्जित प्रतिरक्षा का अध्याय): तैलचित्र,
        Wellcome Collection, Wikimedia Commons के माध्यम से,
        CC~BY~4.0.
  \item \emph{आर्कियोप्टेरिक्स का जीवाश्म} (सामान्य पूर्वजता का अध्याय):
        James L.~Amos का छायाचित्र, Wikimedia Commons, CC0.
\end{itemize}

प्रत्येक छायाचित्र के पूर्ण स्रोत-लिंक और अनुज्ञापत्र के पते पुस्तक के
भंडार में \texttt{images/\allowbreak book2/\allowbreak CREDITS.md} और
\texttt{images/\allowbreak book1/\allowbreak CREDITS.md} में दिए गए हैं।

\bigskip
छायाचित्रों और मौलिक रेखाचित्रों के साथ-साथ, इस खंड के अध्यायों में फैले
लगभग पचास चित्र कृत्रिम बुद्धिमत्ता से बनी चित्रांकन-कृतियाँ हैं, जो
पुस्तक के संपादकों द्वारा लिखे संकेतों से OpenAI की चित्र-जनन प्रणाली से
बनाई गईं और सम्मिलित करने से पहले जैविक शुद्धता के लिए जाँची गईं।
बनाए गए हर चित्र का पूरा संकेत इस किताब के भंडार में, फ़ाइल
\texttt{images/\allowbreak book2/\allowbreak ai/\allowbreak PROMPTS.md} में दर्ज है।
\par\endgroup
\clearpage
